\documentclass[11pt]{article}

\usepackage[T1]{fontenc}
\usepackage[utf8]{inputenc}
\usepackage{lmodern}
\usepackage[margin=1in]{geometry}
\usepackage{amsmath,amssymb,amsthm}
\usepackage{hyperref}
\usepackage{enumitem}
\usepackage{booktabs}
\usepackage{longtable}

\hypersetup{
  colorlinks=true,
  linkcolor=blue,
  citecolor=blue,
  urlcolor=blue,
  pdftitle={OMEGA Survey Memo},
  pdfauthor={Author Name}
}

\newtheorem{conjecture}{Conjecture}
\newtheorem{openquestion}{Open Question}

\title{Replace With A Bounded Survey Title}
\author{Author Name \\ \small OMEGA Research Workflow}
\date{\today}

\begin{document}

\maketitle

\begin{abstract}
Replace this abstract with a concise summary of the survey scope, the number of sources reviewed,
the main structural findings, and the specific open questions that came into focus.
Avoid claiming completeness unless every relevant paper has been checked.
\end{abstract}

\section{Scope and Method}
Describe the exact query surface: which databases were searched, which keyword clusters were used,
what date range was covered, and how borderline relevance was decided.
If AI-assisted screening was used, state the model, prompt strategy, and false-positive rate estimate.

\section{Problem Background}
State the problem or family of problems being surveyed.
Provide the formal statement if one exists, or the strongest currently accepted formulation.

\section{Historical Timeline}
Summarize the key milestones in chronological order.
Use a table when the timeline has more than five entries.

\begin{longtable}{lp{10cm}}
\toprule
Year & Milestone \\
\midrule
Replace & with actual entries \\
\bottomrule
\end{longtable}

\section{Current State of Knowledge}
Organize by approach, technique, or subproblem.
For each cluster, state:
\begin{itemize}[leftmargin=1.5em]
  \item the strongest known result
  \item the evidence class (formal proof, computational bound, heuristic argument)
  \item the gap between current knowledge and the target conjecture
\end{itemize}

\section{AI Amenability Assessment}
Classify the problem along the OMEGA triage dimensions:
\begin{itemize}[leftmargin=1.5em]
  \item Computational tractability (can sub-instances be checked by machine?)
  \item Formalizability (is there a Lean / Coq / Isabelle formulation?)
  \item Search-space structure (is there a well-defined space to explore?)
  \item Prior partial results (are there intermediate claims to build on?)
  \item Data availability (benchmark datasets, sequence databases, tabulations?)
\end{itemize}

\section{Open Questions and Gaps}
List the specific sub-questions that remain open and that are most promising for the next step.

\begin{openquestion}[Replace with a bounded sub-question]
State the sub-question precisely and note its relationship to the main problem.
\end{openquestion}

\section{Recommended Next Steps}
Rank the most productive next actions.
Separate what can be done immediately from what requires new infrastructure.

\section{Limitations of This Survey}
State what was not covered, which sources were inaccessible, and where the AI screening
may have introduced blind spots.

\section*{AI Disclosure}
This survey was prepared with assistance from AI agents for literature discovery,
relevance filtering, timeline assembly, and drafting.
All factual claims require human verification against the cited primary sources.

\section*{Acknowledgments}
Replace or remove this section as appropriate.

\begin{thebibliography}{99}

\bibitem{placeholder}
Replace with actual references.

\end{thebibliography}

\end{document}
